\chapter{Tilt-Table Test}

\section*{What Is This Test?}
A tilt-table test, also called head-up tilt testing, evaluates how blood pressure and heart rate respond when a person moves from lying flat to an upright position. It is mainly used in selected patients with unexplained fainting (syncope), near-fainting, or suspected autonomic dysfunction. The test may reproduce symptoms while heart rhythm, blood pressure, and symptoms are monitored in a controlled setting. It does not identify every cause of fainting and should be interpreted together with the history, examination, and initial cardiovascular evaluation.

\section*{Alternative Names}
Head-Up Tilt Test; Tilt-Table Test; Tilt Testing; Passive Tilt Test; Autonomic Tilt Test; Tilt-Table Study

\section*{Principle}
The patient rests supine on a motorized table secured with safety straps. The table is then tilted, commonly to approximately 60--70 degrees, while continuous or repeated blood-pressure and electrocardiographic monitoring records the cardiovascular response. Some laboratories use a passive phase followed by a medication-provoked phase, such as sublingual nitroglycerin or an intravenous vasoactive agent, when clinically appropriate. A positive response may reproduce symptoms with a characteristic fall in blood pressure, an abnormal heart-rate response, or both; the exact classification depends on the protocol and clinical context.

\section*{Why This Test Is Ordered}
\begin{itemize}
  \item Recurrent unexplained syncope or presyncope after the initial history, examination, orthostatic vital signs, and electrocardiogram have not established a cause
  \item Suspected reflex (vasovagal) syncope when demonstrating a compatible response would change counseling or management
  \item Evaluation of suspected orthostatic hypotension or an abnormal autonomic response when routine measurements are inconclusive
  \item Selected assessment of orthostatic intolerance, including suspected postural orthostatic tachycardia syndrome (POTS), using protocol-specific heart-rate and blood-pressure criteria
  \item Differentiation of syncope from some non-cardiac causes of apparent loss of consciousness in a specialist setting
\end{itemize}

\section*{Primary vs Secondary Test}
Tilt testing is a secondary or confirmatory test, not a routine screening test for every fainting episode. The initial evaluation generally includes a detailed history, physical examination, medication review, orthostatic blood-pressure and heart-rate measurements, and a 12-lead ECG. Ambulatory rhythm monitoring, echocardiography, exercise testing, neurologic assessment, or laboratory testing may be more appropriate when those evaluations suggest another cause.

\section*{How This Test Is Performed}
The patient lies on a motorized table with a footboard and safety straps. An ECG monitors heart rhythm, a blood-pressure cuff or beat-to-beat device monitors blood pressure, and pulse oximetry may be used. After a resting period, the table is raised gradually and the patient remains upright while staff observe symptoms and vital signs. If symptoms become severe or blood pressure falls substantially, the table is returned to the horizontal position. A medication-provoked phase may be added only under the laboratory's protocol and supervision. The study commonly lasts 20--60 minutes, depending on the protocol and recovery time.

\section*{What Preparation Is Needed}
Follow the testing center's instructions about fasting and medications. Many protocols ask the patient to avoid food for several hours and to withhold selected medicines that alter blood pressure or heart rate, but medicines must not be stopped without the ordering clinician's instructions. Wear comfortable clothing, arrange transportation if fainting or medication provocation is possible, and tell staff about pregnancy, heart disease, severe aortic stenosis, recent stroke, or any prior adverse reaction to test medications.

\section*{Contraindications and Precautions}
\begin{itemize}
  \item The test should be deferred or specially planned in patients with unstable cardiovascular disease, severe symptomatic aortic stenosis, uncontrolled arrhythmia, acute myocardial ischemia, severe uncontrolled hypertension, or other conditions in which induced hypotension or syncope would be unsafe. Medication-provoked protocols require additional precautions and are not suitable for everyone. A negative test does not exclude reflex syncope, orthostatic hypotension, or intermittent arrhythmia, and a positive response can occur in people who also have another cause of syncope.
\end{itemize}
\section*{Risks}
The test may cause dizziness, nausea, sweating, palpitations, headache, or temporary loss of consciousness. Blood pressure may fall markedly and the heart rate may slow; staff can lower the table and provide supportive care. Rare risks include prolonged bradycardia, arrhythmia, injury, or complications from a provoking medication. The monitored setting is intended to reduce the risk of injury compared with an unexpected faint at home.

\section*{What Happens After the Test}
The patient remains supine until symptoms and vital signs recover. Temporary fatigue or nausea can occur, especially after medication provocation. The clinician reviews the blood-pressure and heart-rate pattern with the symptoms and other test results. Management may include hydration and positional advice, medication review, physical counter-pressure maneuvers, compression strategies, or referral for further cardiac or autonomic evaluation; the response alone does not determine treatment.

\section*{Understanding Results}

\subsection*{Negative or Normal}
Blood pressure and heart rate remain within the expected range without reproduction of the patient's typical symptoms. This lowers the likelihood of a test-provoked reflex or orthostatic response but does not exclude intermittent disease or other causes of syncope.

\subsection*{Positive or Abnormal}
The patient's typical symptoms are reproduced with a protocol-defined abnormal blood-pressure fall, excessive heart-rate increase, marked bradycardia, or a mixed response. The pattern may support reflex syncope, orthostatic hypotension, or orthostatic tachycardia, but diagnostic thresholds and labels must follow the laboratory protocol and current clinical criteria. An abnormal response should not be interpreted as proof that every prior faint had the same mechanism.

\section*{Follow-up Tests}
\begin{itemize}
  \item \textbf{12-lead ECG and ambulatory ECG monitoring:} To assess for intermittent bradyarrhythmia or tachyarrhythmia when an arrhythmic cause is possible.
  \item \textbf{Orthostatic vital signs or autonomic function testing:} To confirm blood-pressure and heart-rate responses in daily-life positions or to evaluate broader autonomic dysfunction.
  \item \textbf{Echocardiography:} When structural heart disease is suspected from the history, examination, or ECG.
  \item \textbf{Medication review and targeted laboratory tests:} To identify contributors such as volume depletion, anemia, electrolyte disturbance, or drugs that lower blood pressure.
\end{itemize}

\section*{References}
\begin{enumerate}
  \item American Heart Association. Tilt-Table Test. Last reviewed February 26, 2025.
  \item Brignole M, Moya A, de Lange FJ, et al. 2018 ESC Guidelines for the diagnosis and management of syncope. \textit{Eur Heart J}. 2018;39(21):1883--1948.
  \item Shen WK, Sheldon RS, Benditt DG, et al. 2017 ACC/AHA/HRS guideline for the evaluation and management of patients with syncope. \textit{Circulation}. 2017;136:e60--e122.
\end{enumerate}

\clearpage
\section*{Regulatory Status and Authorization Date}
This chapter describes a test category, not a specific commercial product. FDA, CDSCO, EU IVDR, and MHRA approval or registration dates therefore cannot be assigned without the assay manufacturer, product name, instrument, and intended use. For laboratory-developed tests, an individual agency approval date may not exist. See the Preface for the interpretation of regulatory status.

\clearpage
